\documentclass[10pt,a4paper]{amsart}
\usepackage[margin=19mm]{geometry}
\usepackage{amsmath,amssymb,mathtools}
\usepackage[scr=rsfs,cal=euler]{mathalfa}
\usepackage{xfrac}
\usepackage[unicode,hidelinks,bookmarks=false]{hyperref}
\newcommand{\ie}{i.e.\ }
\newcommand{\cf}{cf.\ }
\newcommand{\resp}{respectively\ }
\newcommand{\Subsection}{\subsection}
\providecommand{\nobreakdash}{\nobreak\hbox{-}}
\setlength{\emergencystretch}{2em}
\multlinegap0pt
\usepackage{amssymb}
\usepackage{tikz}
\newtheorem{proposition}{Proposition}
\newcommand{\bal}{\mathrm{b}}
\newcommand{\ext}{\mathrm{e}}
\title{Balance Constants, Majority Cycles,\\ and the Gold Partition Conjecture\\ through Fourteen Elements}
\author{Anish Gupta}
\date{Source: arXiv:2607.23926v2 (CC BY 4.0)}
\begin{document}
\maketitle

\noindent\emph{Source and licence.} Balance Constants, Majority Cycles,\\ and the Gold Partition Conjecture\\ through Fourteen Elements. Version arXiv:2607.23926v2 (CC BY 4.0), distributed by arXiv under the Creative Commons Attribution 4.0 International licence, http://creativecommons.org/licenses/by/4.0/, as recorded in the arXiv licence metadata for this exact version and shown on the abstract page https://arxiv.org/abs/2607.23926v2. Full licence notice: \texttt{licenses/arxiv-2607.23926-CC-BY-4.0.txt}.\par\smallskip

\noindent\textit{Editorial scope.} Counted unit: the article's proposition prop:equality (source0.tex lines 182--191). The article's complete original proof of it is delivered with it (source0.tex lines 585--594, one \texttt{proof} environment of 10 lines, which the article places away from the statement). No mathematics is written, repaired or re-derived: the statement and the proof are the article's own lines, in the article's own notation, and no retained line is substituted or re-flowed.  Retained uncounted as the source-local context the proof needs: lines 362--364.  Nothing was omitted from any retained span, so the package carries no omission record.  No retained line contains a citation, so the wrapper carries no bibliography and no bibliography entry.  No retained line contains a figure environment, an image-inclusion command or drawing code, so no image is delivered and the package declares no asset dependency.  Editorial formatting only, in addition: an AMS \texttt{amsart} preamble that restates the article's own declarations and macros used by the retained lines, namely the environments \texttt{proposition} on separate counters, together with the standard packages those lines need.  The retained environments print in this document as Proposition 1 in order of appearance, whereas the article prints the same items with the numbering of its own full document.  The complete native source of the article as distributed in the arXiv e-print for this exact version is bundled unchanged as \texttt{sources/206078/source0.tex}.
\begin{equation}\label{eq:sumbal}
  \bal(P)=\max_i\bal(P_i),
\end{equation}

\begin{proposition}\label{prop:equality}
Let $T$ be the three-element poset consisting of a two-element chain and an
isolated point.  Every ordinal sum of singletons and copies of $T$, other than
the chain, has balance constant exactly $1/3$, and the number of such posets
on $n$ elements is $a(n)-1$, where
\[
  a(n)=a(n-1)+a(n-3),\qquad a(0)=a(1)=a(2)=1.
\]
No other poset on at most $14$ elements has balance constant $1/3$.
\end{proposition}

\begin{proof}[Proof of Proposition~\ref{prop:equality}]
The poset $T$ has $\ext(T)=3$ and $\bal(T)=1/3$.  By \eqref{eq:sumbal}, an
ordinal sum of singletons and copies of $T$ has balance constant $1/3$ as soon
as at least one summand is a copy of $T$; if all summands are singletons the
sum is a chain and has no incomparable pair.  Decomposition into
ordinal-indecomposable summands is unique, so these posets are in bijection
with the compositions of $n$ into parts $1$ and $3$.  Splitting on the first
part gives $a(n)=a(n-1)+a(n-3)$, and discarding the all-singleton composition
gives $a(n)-1$.
\end{proof}

\end{document}
